\documentclass{beamer}

\usepackage{beamerthemeSingapore}
\usepackage[brazil]{babel}
\usepackage[T1]{fontenc}
\usepackage[utf8]{inputenc}
\usepackage{graphicx}
\usepackage{url}
\usepackage{xcolor}
\usepackage{booktabs}
\usepackage{ragged2e}
\usepackage{etoolbox}
\usepackage{amsmath}
\usepackage{minted}
\usepackage{tikz}
\usetikzlibrary{positioning, arrows.meta, shapes.geometric}

\newcommand{\tcb}[1]{\textcolor{blue}{#1}}
\newcommand{\tcr}[1]{\textcolor{red}{#1}}
\newcommand{\tcg}[1]{\textcolor{green!50!black}{#1}}

\tikzset{
  treenode/.style = {align=center, inner sep=0pt, text centered,
    font=\sffamily\small},
  arn_b/.style = {treenode, circle, black, draw=black,
    fill=white, text width=1.7em},
  arn_d/.style = {treenode, circle, white, draw=black,
    fill=blue!60, text width=1.7em},
  arn_r/.style = {treenode, circle, white, draw=black,
    fill=red!70!black, text width=1.7em},
  arn_t/.style = {treenode, regular polygon, regular polygon sides=3,
    draw=black, fill=black!8, minimum size=1.1cm, inner sep=0pt,
    font=\sffamily\footnotesize}
}

\setminted{
  fontsize=\footnotesize,
  breaklines=true,
  breakanywhere=true,
  frame=lines,
  framesep=2mm,
  baselinestretch=1.05,
  tabsize=2
}

\apptocmd{\frame}{}{\justifying}{}

\setbeamertemplate{navigation symbols}{}

\title{06 - Árvores binárias de busca e AVL}
\author{Prof. Alex Kutzke}
\date{Setembro de 2026}

\begin{document}

% ==========================================
\begin{frame}
  \begin{center}
    \large{Setor de Educação Profissional e Tecnológica}\\
    Tecnologia em Análise e Desenvolvimento de Sistemas\\
    \vspace{1cm}
    \small{DS143 - Estruturas de Dados II}
  \end{center}
  \titlepage
\end{frame}

% ==========================================
\begin{frame}{Objetivos da aula}
  Ao final desta aula você deve ser capaz de:
  \begin{enumerate}
    \item Enunciar a propriedade da árvore binária de busca e verificar se uma
          árvore dada a satisfaz;
    \item Implementar busca, inserção, mínimo e máximo em Go, e justificar por
          que o custo das quatro operações é a altura da árvore;
    \item Executar à mão a remoção nos três casos;
    \item Explicar por que a ordem de inserção determina a forma da árvore;
    \item Calcular o fator de balanceamento e identificar qual dos quatro casos
          de desbalanceamento ocorreu;
    \item Aplicar rotações simples e duplas;
    \item Comparar os custos de ABB e AVL no melhor caso, no caso médio e no
          pior caso.
  \end{enumerate}

  \vspace{0.3cm}
  \footnotesize Este conteúdo cai na Prova 1.
\end{frame}

% ==========================================
\begin{frame}{Roteiro}
  \tableofcontents
\end{frame}

% ==========================================
\section{Motivação e definição}

\begin{frame}{Guardar um conjunto de chaves}
  Duas operações: inserir uma chave nova, e responder se uma chave está no
  conjunto.

  \vspace{0.5cm}
  \begin{center}
    \begin{tabular}{lll}
      \toprule
      \textbf{Estrutura} & \textbf{Inserção} & \textbf{Busca} \\
      \midrule
      Slice sem ordem & $O(1)$ amortizado & $\Theta(n)$ \\
      Slice ordenado  & $\Theta(n)$ & $\Theta(\log n)$ busca binária \\
      \bottomrule
    \end{tabular}
  \end{center}

  \vspace{0.5cm}
  No slice ordenado, inserir no meio obriga a empurrar para a direita, uma a
  uma, todas as chaves que vêm depois da chave nova.
\end{frame}

\begin{frame}[fragile]{A medição: 200 mil inserções, 10 mil buscas}
\begin{minted}{text}
$ go run motivacao.go
200000 chaves inseridas, 10000 buscas

estrutura        | inserção | busca
-----------------+----------+---------
slice sem ordem  |      3ms | 325ms
slice ordenado   |   2.334s | 1ms
árvore de busca  |     54ms | 2ms
\end{minted}

  \vspace{0.2cm}
  A árvore de busca fica perto do melhor das duas colunas. Ela não precisa
  manter as chaves contíguas na memória: a ordem está na forma da árvore, e
  mudar a forma custa alterar um ponteiro.
\end{frame}

\begin{frame}[fragile]{A busca que a aula anterior deixou em aberto}
\begin{minted}{go}
func Pertence(a *No, procurado int) bool {
	if a == nil {
		return false
	}
	if a.Info == procurado {
		return true
	}
	return Pertence(a.Esq, procurado) ||
		Pertence(a.Dir, procurado)
}
\end{minted}

  \vspace{0.2cm}
  Visita todos os nós no pior caso: nada na árvore diz de que lado continuar.

  \vspace{0.2cm}
  Mas o in-ordem da árvore de exemplo saiu crescente:
  \texttt{10 20 30 35 40 45 50 90 95}. Aquela árvore tinha sido montada com os
  menores à esquerda e os maiores à direita.
\end{frame}

\begin{frame}{Árvore binária de busca}
  \begin{columns}
    \begin{column}{0.42\textwidth}
      \centering
      \begin{tikzpicture}[-, level distance=0.9cm,
          level 1/.style={sibling distance=2.6cm},
          level 2/.style={sibling distance=1.3cm},
          level 3/.style={sibling distance=1cm}]
        \node [arn_b] {50}
        child { node [arn_b] {30}
            child { node [arn_b] {20}
                child { node [arn_b] {10} }
                child [missing]
              }
            child { node [arn_b] {40}
                child { node [arn_b] {35} }
                child { node [arn_b] {45} }
              }
          }
        child { node [arn_b] {90}
            child [missing]
            child { node [arn_b] {95} }
          };
      \end{tikzpicture}
    \end{column}
    \begin{column}{0.58\textwidth}
      Árvore binária em que, para \tcr{todo} nó $x$:
      \begin{itemize}
        \item as chaves da subárvore esquerda de $x$ são menores que a chave de
              $x$;
        \item as chaves da subárvore direita de $x$ são maiores que a chave de
              $x$.
      \end{itemize}

      \vspace{0.3cm}
      Nesta aula, chaves distintas.
    \end{column}
  \end{columns}
\end{frame}

\begin{frame}{A exigência vale para todos os nós}
  \begin{columns}
    \begin{column}{0.4\textwidth}
      \centering
      \begin{tikzpicture}[-, level distance=1cm,
          level 1/.style={sibling distance=2.4cm},
          level 2/.style={sibling distance=1.4cm}]
        \node [arn_b] {50}
        child { node [arn_b] {30}
            child [missing]
            child { node [arn_r] {60} }
          }
        child { node [arn_b] {90} };
      \end{tikzpicture}
    \end{column}
    \begin{column}{0.6\textwidth}
      Cada nó está do lado certo em relação ao \textbf{seu pai}, e ainda assim
      esta árvore \tcr{não é uma ABB}: o 60 está na subárvore esquerda do 50.

      \vspace{0.3cm}
      O erro clássico na função que verifica a propriedade é comparar cada nó
      apenas com os seus dois filhos. Esse teste aceita a árvore ao lado.
    \end{column}
  \end{columns}
\end{frame}

\begin{frame}{Duas consequências da definição}
  \begin{block}{O percurso in-ordem devolve as chaves em ordem crescente}
    A garantia que a árvore da aula anterior tinha por construção agora é
    obrigatória. É o teste mais rápido para saber se uma operação preservou a
    propriedade de busca.
  \end{block}

  \begin{block}{A busca nunca precisa voltar atrás}
    Em cada nó, a comparação descarta uma subárvore inteira. O caminho
    percorrido é único, e tem no máximo a altura da árvore mais um nó.
  \end{block}
\end{frame}

% ==========================================
\section{Operações}

\begin{frame}[fragile]{Busca}
\begin{minted}{go}
func Busca(a *No, procurado int) bool {
	if a == nil {
		return false
	}
	if procurado < a.Info {
		return Busca(a.Esq, procurado)
	}
	if procurado > a.Info {
		return Busca(a.Dir, procurado)
	}
	return true
}
\end{minted}

  \vspace{0.2cm}
  Três resultados na comparação, e cada um decide o passo seguinte.
\end{frame}

\begin{frame}[fragile]{A mesma busca, sem recursão}
\begin{minted}{go}
func BuscaIterativa(a *No, procurado int) bool {
	for a != nil {
		if procurado < a.Info {
			a = a.Esq
		} else if procurado > a.Info {
			a = a.Dir
		} else {
			return true
		}
	}
	return false
}
\end{minted}

  \vspace{0.2cm}
  As chamadas recursivas estavam na posição final e nenhuma precisava do
  resultado da outra. A \texttt{Pertence} do slide anterior precisa das duas
  subárvores, e não vira laço com a mesma facilidade.
\end{frame}

\begin{frame}{O caminho da busca por 35}
  \begin{center}
    \begin{tikzpicture}[-, level distance=1.1cm,
        level 1/.style={sibling distance=4cm},
        level 2/.style={sibling distance=2cm},
        level 3/.style={sibling distance=1.4cm}]
      \node [arn_d] {50}
      child { node [arn_d] {30}
          child { node [arn_b] {20}
              child { node [arn_b] {10} }
              child [missing]
            }
          child { node [arn_d] {40}
              child { node [arn_d] {35} }
              child { node [arn_b] {45} }
            }
        }
      child { node [arn_b] {90}
          child [missing]
          child { node [arn_b] {95} }
        };
    \end{tikzpicture}
  \end{center}

  \vspace{0.2cm}
  Quatro comparações, quatro nós visitados de nove. Em cada nó, metade do que
  restava é descartada sem ser olhada.
\end{frame}

\begin{frame}[fragile]{Inserção}
\begin{minted}{go}
func Insere(a *No, v int) *No {
	if a == nil {
		return &No{Info: v}
	}
	if v < a.Info {
		a.Esq = Insere(a.Esq, v)
	} else if v > a.Info {
		a.Dir = Insere(a.Dir, v)
	}
	return a
}
\end{minted}

  \vspace{0.2cm}
  O nó novo entra como \textbf{folha}, na posição em que a busca por aquela
  chave terminaria. Quem chama escreve \texttt{a = Insere(a, v)}, o que cobre
  também a árvore vazia.
\end{frame}

\begin{frame}{Onde entra o 37}
  \begin{columns}
    \begin{column}{0.55\textwidth}
      \centering
      \begin{tikzpicture}[-, level distance=0.95cm,
          level 1/.style={sibling distance=2.8cm},
          level 2/.style={sibling distance=1.5cm},
          level 3/.style={sibling distance=1.1cm}]
        \node [arn_b] {50}
        child { node [arn_b] {30}
            child { node [arn_b] {20}
                child { node [arn_b] {10} }
                child [missing]
              }
            child { node [arn_b] {40}
                child { node [arn_b] {35} }
                child { node [arn_b] {45} }
              }
          }
        child { node [arn_b] {90}
            child [missing]
            child { node [arn_b] {95} }
          };
      \end{tikzpicture}
    \end{column}
    \begin{column}{0.45\textwidth}
      \begin{block}{}
        \texttt{a = Insere(a, 37)}

        \vspace{0.3cm}
        De qual nó o 37 vira filho, e de que lado?

        \vspace{0.2cm}
        Quantas comparações até chegar lá?
      \end{block}
    \end{column}
  \end{columns}
\end{frame}

\begin{frame}{Onde entra o 37}
  \begin{columns}
    \begin{column}{0.55\textwidth}
      \centering
      \begin{tikzpicture}[-, level distance=0.95cm,
          level 1/.style={sibling distance=2.8cm},
          level 2/.style={sibling distance=1.5cm},
          level 3/.style={sibling distance=1.1cm}]
        \node [arn_d] {50}
        child { node [arn_d] {30}
            child { node [arn_b] {20}
                child { node [arn_b] {10} }
                child [missing]
              }
            child { node [arn_d] {40}
                child { node [arn_d] {35}
                    child [missing]
                    child { node [arn_r] {37} }
                  }
                child { node [arn_b] {45} }
              }
          }
        child { node [arn_b] {90}
            child [missing]
            child { node [arn_b] {95} }
          };
      \end{tikzpicture}
    \end{column}
    \begin{column}{0.45\textwidth}
      $37 < 50$: esquerda.\\
      $37 > 30$: direita.\\
      $37 < 40$: esquerda.\\
      $37 > 35$: direita, e ali está o \texttt{nil}.

      \vspace{0.4cm}
      A linha \texttt{a.Esq = Insere(a.Esq, v)} reatribui o ponteiro em todos os
      nós do caminho, embora só um deles mude. É o preço de não guardar o
      ponteiro para o pai.
    \end{column}
  \end{columns}
\end{frame}

\begin{frame}[fragile]{Mínimo, máximo e listagem em ordem}
\begin{minted}{go}
func Minimo(a *No) (int, bool) {
	if a == nil {
		return 0, false
	}
	for a.Esq != nil {
		a = a.Esq
	}
	return a.Info, true
}
\end{minted}

  \vspace{0.2cm}
  O menor valor está no nó mais à esquerda, o maior no mais à direita. Nenhuma
  comparação de chaves é necessária.

  \vspace{0.2cm}
  O segundo valor devolvido diz se a resposta existe. A árvore vazia não tem
  mínimo, e devolver 0 sozinho a confundiria com uma árvore que contém o 0.
\end{frame}

\begin{frame}{Remoção: dois casos fáceis}
  \begin{columns}[t]
    \begin{column}{0.5\textwidth}
      \centering
      \textbf{Caso 1: folha}\\[0.2cm]
      \begin{tikzpicture}[-, level distance=0.9cm,
          level 1/.style={sibling distance=1.6cm}]
        \node [arn_b] {40}
        child { node [arn_b] {35} }
        child { node [arn_r] {45} };
      \end{tikzpicture}\\[0.2cm]
      O nó some, e o pai passa a apontar para \texttt{nil}.
    \end{column}
    \begin{column}{0.5\textwidth}
      \centering
      \textbf{Caso 2: uma subárvore só}\\[0.2cm]
      \begin{tikzpicture}[-, level distance=0.9cm,
          level 1/.style={sibling distance=1.6cm}]
        \node [arn_b] {30}
        child { node [arn_r] {20}
            child { node [arn_b] {10} }
            child [missing]
          }
        child { node [arn_b] {40} };
      \end{tikzpicture}\\[0.2cm]
      O filho sobe para o lugar do nó removido, com toda a sua subárvore.
    \end{column}
  \end{columns}

  \vspace{0.4cm}
  No caso 2 a propriedade continua valendo: tudo o que subiu já estava do lado
  certo em relação ao avô.
\end{frame}

\begin{frame}{Remoção: duas subárvores}
  \begin{columns}
    \begin{column}{0.45\textwidth}
      \centering
      \begin{tikzpicture}[-, level distance=0.95cm,
          level 1/.style={sibling distance=2.6cm},
          level 2/.style={sibling distance=1.4cm},
          level 3/.style={sibling distance=1.1cm}]
        \node [arn_r] {50}
        child { node [arn_b] {30}
            child { node [arn_b] {20}
                child { node [arn_b] {10} }
                child [missing]
              }
            child { node [arn_b] {40}
                child { node [arn_b] {35} }
                child { node [arn_b] {45} }
              }
          }
        child { node [arn_b] {90}
            child [missing]
            child { node [arn_b] {95} }
          };
      \end{tikzpicture}
    \end{column}
    \begin{column}{0.55\textwidth}
      \begin{block}{}
        \texttt{a = Remove(a, 50)}

        \vspace{0.3cm}
        Nenhum dos dois filhos pode simplesmente subir.

        \vspace{0.3cm}
        Que chave da árvore pode ocupar o lugar do 50? Quantas servem?
      \end{block}
    \end{column}
  \end{columns}
\end{frame}

\begin{frame}{Remoção: duas subárvores}
  \begin{columns}
    \begin{column}{0.45\textwidth}
      \centering
      \begin{tikzpicture}[-, level distance=0.95cm,
          level 1/.style={sibling distance=2.6cm},
          level 2/.style={sibling distance=1.4cm},
          level 3/.style={sibling distance=1.1cm}]
        \node [arn_r] {50}
        child { node [arn_b] {30}
            child { node [arn_b] {20}
                child { node [arn_b] {10} }
                child [missing]
              }
            child { node [arn_b] {40}
                child { node [arn_b] {35} }
                child { node [arn_d] {45} }
              }
          }
        child { node [arn_b] {90}
            child [missing]
            child { node [arn_b] {95} }
          };
      \end{tikzpicture}
    \end{column}
    \begin{column}{0.55\textwidth}
      O lugar do 50 comporta uma chave maior que tudo à esquerda e menor que
      tudo à direita. Existem duas assim:
      \begin{itemize}
        \item o \tcb{antecessor}, maior valor da subárvore esquerda (45);
        \item o \textbf{sucessor}, menor valor da subárvore direita (90).
      \end{itemize}

      \vspace{0.3cm}
      Copiamos uma das duas para o nó e removemos o nó de onde ela veio.

      \vspace{0.3cm}
      O antecessor não tem filho à direita, então essa segunda remoção cai no
      caso 1 ou no caso 2. \tcr{Nunca no caso 3.}
    \end{column}
  \end{columns}
\end{frame}

\begin{frame}[fragile]{Remoção em Go}
\begin{minted}{go}
	// achou o nó: três casos
	if a.Esq == nil && a.Dir == nil {
		return nil
	}
	if a.Esq == nil {
		return a.Dir
	}
	if a.Dir == nil {
		return a.Esq
	}

	antecessor := a.Esq
	for antecessor.Dir != nil {
		antecessor = antecessor.Dir
	}
	a.Info = antecessor.Info
	a.Esq = Remove(a.Esq, antecessor.Info)
	return a
\end{minted}
\end{frame}

\begin{frame}[fragile]{Remoção: as três saídas}
\begin{minted}{text}
Remove(a, 45), caso 1: folha
in-ordem: 10 20 30 35 40 50 90 95

Remove(a, 20), caso 2: uma subárvore só
in-ordem: 10 30 35 40 50 90 95

Remove(a, 50), caso 3: duas subárvores (a raiz)
        95
    90
40
        35
    30
        10
in-ordem: 10 30 35 40 90 95
\end{minted}

  \vspace{0.2cm}
  O 40, maior valor da subárvore esquerda, tomou o lugar do 50 na raiz. O
  in-ordem continua crescente em todas as saídas.
\end{frame}

% ==========================================
\section{Forma e desempenho}

\begin{frame}{A ordem de inserção determina a forma}
  \begin{columns}
    \begin{column}{0.5\textwidth}
      \centering
      \footnotesize \texttt{50 30 90 20 40 95 10 35 45}\\[0.2cm]
      \begin{tikzpicture}[-, level distance=0.8cm,
          level 1/.style={sibling distance=2.2cm},
          level 2/.style={sibling distance=1.1cm},
          level 3/.style={sibling distance=0.9cm},
          every node/.style={font=\sffamily\scriptsize}]
        \node [arn_b] {50}
        child { node [arn_b] {30}
            child { node [arn_b] {20}
                child { node [arn_b] {10} }
                child [missing]
              }
            child { node [arn_b] {40}
                child { node [arn_b] {35} }
                child { node [arn_b] {45} }
              }
          }
        child { node [arn_b] {90}
            child [missing]
            child { node [arn_b] {95} }
          };
      \end{tikzpicture}\\[0.2cm]
      altura 3
    \end{column}
    \begin{column}{0.5\textwidth}
      \centering
      \footnotesize \texttt{10 20 30 35 40 45 50 90 95}\\[0.2cm]
      \begin{tikzpicture}[-, level distance=0.55cm,
          level/.style={sibling distance=0.7cm},
          every node/.style={font=\sffamily\scriptsize}]
        \node [arn_b] {10}
        child [missing]
        child { node [arn_b] {20}
            child [missing]
            child { node [arn_b] {30}
                child [missing]
                child { node [arn_b] {35}
                    child [missing]
                    child { node [arn_b] {40}
                        child [missing]
                        child { node [arn_b] {45}
                            child [missing]
                            child { node [arn_b] {50}
                                child [missing]
                                child { node [arn_b] {90}
                                    child [missing]
                                    child { node [arn_b] {95} }
                                  }
                              }
                          }
                      }
                  }
              }
          };
      \end{tikzpicture}\\[0.2cm]
      altura 8
    \end{column}
  \end{columns}

  \vspace{0.3cm}
  \centering
  Mesmos nove valores. A segunda é uma lista encadeada com passos extras.
\end{frame}

\begin{frame}{Entrada ordenada é comum}
  Chaves lidas de um arquivo já ordenado, códigos sequenciais e carimbos de
  tempo produzem a árvore degenerada, justamente nos casos em que a árvore
  seria mais útil.

  \vspace{0.5cm}
  Entre os dois extremos está o caso médio, com as chaves chegando em ordem
  aleatória:

  \vspace{0.3cm}
  \begin{center}
    \footnotesize
    \begin{tabular}{rrrrr}
      \toprule
      $n$ & mínima & aleatória & degenerada & comparações médias \\
      \midrule
      1.000 & 9 & 21,1 & 999 & 12,0 \\
      10.000 & 13 & 30,6 & 9.999 & 16,6 \\
      100.000 & 16 & 40,6 & 99.999 & 21,0 \\
      1.000.000 & 19 & 46,8 & 999.999 & 25,3 \\
      \bottomrule
    \end{tabular}
  \end{center}

  \vspace{0.3cm}
  O número médio de comparações de uma busca bem sucedida tende a
  $2 \ln n \approx 1{,}39 \log_2 n$, cerca de 39\% acima da árvore
  perfeitamente balanceada.
\end{frame}

\begin{frame}{O custo de todas as operações é a altura}
  Busca, inserção, remoção, mínimo e máximo percorrem um único caminho da raiz
  para baixo: $O(h)$.

  \vspace{0.4cm}
  \begin{center}
    \footnotesize
    \begin{tabular}{lll}
      \toprule
      \textbf{Situação} & \textbf{Altura} & \textbf{Custo} \\
      \midrule
      Árvore cheia (melhor caso) & $\log_2(n+1) - 1$ & $\Theta(\log n)$ \\
      Ordem aleatória (caso médio) & $\Theta(\log n)$ & $\Theta(\log n)$ \\
      Ordem crescente (pior caso) & $n - 1$ & $\Theta(n)$ \\
      \bottomrule
    \end{tabular}
  \end{center}

  \vspace{0.5cm}
  O caso médio supõe que a ordem de chegada das chaves seja aleatória, e quem
  escreve a estrutura não controla essa ordem.
\end{frame}

% ==========================================
\section{Balanceamento AVL}

\begin{frame}{O critério AVL}
  Manter a árvore perfeitamente balanceada custaria caro: uma única inserção
  pode obrigar a reconstruir quase tudo. Adelson-Velsky e Landis, em 1962,
  exigiram menos.

  \vspace{0.4cm}
  \begin{block}{Árvore AVL}
    Árvore binária de busca em que, para todo nó, as alturas das duas
    subárvores diferem em no máximo 1.
  \end{block}

  \vspace{0.4cm}
  \textbf{Fator de balanceamento} de um nó:
  $$\mathrm{fb}(x) = \mathrm{altura}(x.\mathrm{Esq}) - \mathrm{altura}(x.\mathrm{Dir})$$

  Em uma AVL, todo nó tem fb igual a $-1$, 0 ou $+1$. Com fb $= +2$, o nó está
  pesado à esquerda; com fb $= -2$, à direita. A árvore vazia tem altura $-1$.
\end{frame}

\begin{frame}{Fator de balanceamento na árvore de exemplo}
  \begin{center}
    \begin{tikzpicture}[-, level distance=1.1cm,
        level 1/.style={sibling distance=4cm},
        level 2/.style={sibling distance=2cm},
        level 3/.style={sibling distance=1.4cm}]
      \node [arn_b] {50}
      child { node [arn_b] {30}
          child { node [arn_b] {20}
              child { node [arn_b] {10} }
              child [missing]
            }
          child { node [arn_b] {40}
              child { node [arn_b] {35} }
              child { node [arn_b] {45} }
            }
        }
      child { node [arn_b] {90}
          child [missing]
          child { node [arn_b] {95} }
        };
    \end{tikzpicture}
  \end{center}

  \vspace{0.2cm}
  Qual é o fator de balanceamento de cada um dos cinco nós internos? A árvore é
  AVL?
\end{frame}

\begin{frame}{Fator de balanceamento na árvore de exemplo}
  \begin{center}
    \begin{tikzpicture}[-, level distance=1.1cm,
        level 1/.style={sibling distance=4cm},
        level 2/.style={sibling distance=2cm},
        level 3/.style={sibling distance=1.4cm}]
      \node [arn_b, label=right:{\scriptsize fb $+1$}] {50}
      child { node [arn_b, label=left:{\scriptsize fb $0$}] {30}
          child { node [arn_b, label=left:{\scriptsize fb $+1$}] {20}
              child { node [arn_b] {10} }
              child [missing]
            }
          child { node [arn_b, label=right:{\scriptsize fb $0$}] {40}
              child { node [arn_b] {35} }
              child { node [arn_b] {45} }
            }
        }
      child { node [arn_b, label=right:{\scriptsize fb $-1$}] {90}
          child [missing]
          child { node [arn_b] {95} }
        };
    \end{tikzpicture}
  \end{center}

  \vspace{0.2cm}
  Todos os fatores estão em $\{-1, 0, +1\}$: a árvore de exemplo é AVL. A
  degenerada do slide anterior não é, a raiz dela tem fb $= -8$.
\end{frame}

\begin{frame}[fragile]{A altura guardada no nó}
\begin{minted}[fontsize=\scriptsize]{go}
type No struct {
	Info int
	Alt  int
	Esq  *No
	Dir  *No
}

func altura(a *No) int {
	if a == nil {
		return -1
	}
	return a.Alt
}

func atualizaAltura(a *No) {
	a.Alt = 1 + maior(altura(a.Esq), altura(a.Dir))
}
\end{minted}

  \vspace{0.2cm}
  Recalcular a altura a cada verificação custaria $\Theta(n)$. Com o campo
  \texttt{Alt}, calcular o fator de balanceamento custa tempo constante.
\end{frame}

\begin{frame}{Por que o critério fraco basta}
  Seja $N(h)$ o \tcr{menor} número de nós de uma AVL de altura $h$. A árvore
  mais esguia possível tem uma raiz, uma subárvore de altura $h-1$ e outra tão
  pequena quanto o critério permite, de altura $h-2$:

  $$N(h) = N(h-1) + N(h-2) + 1, \quad N(0) = 1, \quad N(1) = 2$$

  \vspace{0.2cm}
  \begin{center}
    \footnotesize
    \begin{tabular}{lrrrrrrrrrrr}
      \toprule
      $h$ & 0 & 1 & 2 & 3 & 4 & 5 & 6 & 7 & 8 & 9 & 10 \\
      $N(h)$ & 1 & 2 & 4 & 7 & 12 & 20 & 33 & 54 & 88 & 143 & 232 \\
      \bottomrule
    \end{tabular}
  \end{center}

  \vspace{0.3cm}
  É a recorrência de Fibonacci deslocada, $N(h) = F(h+3) - 1$, e cresce
  exponencialmente com $h$. Invertendo:

  $$h < 1{,}44 \log_2 (n + 2)$$

  Uma AVL com 10 mil chaves tem altura no máximo 17: a menor AVL de altura 18
  já precisaria de 10.945 nós.
\end{frame}

% ==========================================
\section{Rotações}

\begin{frame}{Rotação à direita}
  \begin{columns}
    \begin{column}{0.45\textwidth}
      \centering
      \begin{tikzpicture}[-, level distance=1cm,
          level 1/.style={sibling distance=2cm},
          level 2/.style={sibling distance=1.4cm}]
        \node [arn_b] {a}
        child { node [arn_b] {b}
            child { node [arn_t] {x} }
            child { node [arn_t] {y} }
          }
        child { node [arn_t] {z} };
      \end{tikzpicture}
    \end{column}
    \begin{column}{0.55\textwidth}
      \begin{block}{}
        A subárvore está pesada à esquerda. Queremos que \textbf{b} suba para a
        raiz e \textbf{a} desça para a direita.

        \vspace{0.3cm}
        Onde ficam $x$, $y$ e $z$ depois disso?
      \end{block}
    \end{column}
  \end{columns}

  \vspace{0.4cm}
  Só uma das três subárvores muda de pai. Qual, e por quê?
\end{frame}

\begin{frame}{Rotação à direita}
  \begin{columns}
    \begin{column}{0.45\textwidth}
      \centering
      \begin{tikzpicture}[-, level distance=1cm,
          level 1/.style={sibling distance=2cm},
          level 2/.style={sibling distance=1.4cm}]
        \node [arn_b] {a}
        child { node [arn_b] {b}
            child { node [arn_t] {x} }
            child { node [arn_t] {y} }
          }
        child { node [arn_t] {z} };
      \end{tikzpicture}
    \end{column}
    \begin{column}{0.1\textwidth}
      \centering
      \scriptsize rotação\\ à direita
    \end{column}
    \begin{column}{0.45\textwidth}
      \centering
      \begin{tikzpicture}[-, level distance=1cm,
          level 1/.style={sibling distance=2cm},
          level 2/.style={sibling distance=1.4cm}]
        \node [arn_b] {b}
        child { node [arn_t] {x} }
        child { node [arn_b] {a}
            child { node [arn_t] {y} }
            child { node [arn_t] {z} }
          };
      \end{tikzpicture}
    \end{column}
  \end{columns}

  \vspace{0.5cm}
  In-ordem da árvore da esquerda: $x$, b, $y$, a, $z$.\\
  In-ordem da árvore da direita: $x$, b, $y$, a, $z$.

  \vspace{0.3cm}
  As duas sequências são iguais, e por isso a rotação nunca viola a propriedade
  de busca. O que muda é a altura: b sobe um nível com $x$, a desce um nível
  com $z$.
\end{frame}

\begin{frame}[fragile]{Rotação à direita em Go}
\begin{minted}{go}
func RotacaoDireita(a *No) *No {
	b := a.Esq
	a.Esq = b.Dir
	b.Dir = a

	atualizaAltura(a)
	atualizaAltura(b)
	return b
}
\end{minted}

  \vspace{0.2cm}
  Três atribuições de ponteiro, e a ordem entre elas importa: \texttt{b.Dir} é
  lido para dentro de \texttt{a.Esq} antes de ser sobrescrito com \texttt{a}.

  \vspace{0.2cm}
  As alturas são atualizadas de baixo para cima. A função devolve a nova raiz da
  subárvore, e quem chamou reatribui o ponteiro.
  \texttt{RotacaoEsquerda} é a imagem espelhada.
\end{frame}

\begin{frame}{Os quatro casos}
  Depois de uma inserção, o primeiro nó desbalanceado na volta da recursão tem
  fb $= +2$ ou fb $= -2$. A rotação a aplicar depende também do lado em que o
  filho pesado cresceu.

  \vspace{0.3cm}
  \begin{center}
    \footnotesize
    \begin{tabular}{llll}
      \toprule
      \textbf{Caso} & \textbf{fb do nó} & \textbf{fb do filho} & \textbf{Correção} \\
      \midrule
      Esquerda-esquerda & $+2$ & $+1$ & simples à direita \\
      Direita-direita & $-2$ & $-1$ & simples à esquerda \\
      Esquerda-direita & $+2$ & $-1$ & esquerda no filho, direita no nó \\
      Direita-esquerda & $-2$ & $+1$ & direita no filho, esquerda no nó \\
      \bottomrule
    \end{tabular}
  \end{center}

  \vspace{0.4cm}
  Nos dois primeiros casos o nó e o filho pendem para o mesmo lado, e uma
  rotação resolve. Nos dois últimos o desequilíbrio faz um zigue-zague, e uma
  rotação simples apenas o inverteria.
\end{frame}

\begin{frame}{Caso esquerda-esquerda: inserir 30, 20, 10}
  \begin{columns}
    \begin{column}{0.45\textwidth}
      \centering
      \begin{tikzpicture}[-, level distance=1cm,
          level 1/.style={sibling distance=1.6cm}]
        \node [arn_b, label=right:{\scriptsize fb $+2$}] {30}
        child { node [arn_b, label=right:{\scriptsize fb $+1$}] {20}
            child { node [arn_r] {10} }
            child [missing]
          }
        child [missing];
      \end{tikzpicture}
    \end{column}
    \begin{column}{0.1\textwidth}
      \centering
      \scriptsize rotação\\ à direita
    \end{column}
    \begin{column}{0.45\textwidth}
      \centering
      \begin{tikzpicture}[-, level distance=1cm,
          level 1/.style={sibling distance=1.8cm}]
        \node [arn_b, label=right:{\scriptsize fb $0$}] {20}
        child { node [arn_b] {10} }
        child { node [arn_b] {30} };
      \end{tikzpicture}
    \end{column}
  \end{columns}

  \vspace{0.6cm}
  \centering
  Uma rotação simples à direita no 30. Altura 2 antes, altura 1 depois.
\end{frame}

\begin{frame}{Caso esquerda-direita: inserir 30, 10, 20}
  \begin{columns}
    \begin{column}{0.4\textwidth}
      \centering
      \begin{tikzpicture}[-, level distance=1cm,
          level 1/.style={sibling distance=1.6cm}]
        \node [arn_b, label=right:{\scriptsize fb $+2$}] {30}
        child { node [arn_b] {10}
            child [missing]
            child { node [arn_r] {20} }
          }
        child [missing];
      \end{tikzpicture}
    \end{column}
    \begin{column}{0.6\textwidth}
      \begin{block}{}
        O 30 está pesado à esquerda, como no caso anterior.

        \vspace{0.3cm}
        Uma rotação simples à direita no 30 resolve?
      \end{block}
    \end{column}
  \end{columns}

  \vspace{0.5cm}
  \centering
  Compare o fator do 30 com o do filho 10. No caso anterior os dois tinham o
  mesmo sinal.
\end{frame}

\begin{frame}{Caso esquerda-direita: inserir 30, 10, 20}
  \begin{columns}
    \begin{column}{0.24\textwidth}
      \centering
      \begin{tikzpicture}[-, level distance=0.9cm,
          level 1/.style={sibling distance=1.3cm}]
        \node [arn_b, label=right:{\scriptsize fb $+2$}] {30}
        child { node [arn_b] {10}
            child [missing]
            child { node [arn_r] {20} }
          }
        child [missing];
      \end{tikzpicture}
    \end{column}
    \begin{column}{0.13\textwidth}
      \centering \scriptsize à esquerda\\ no 10
    \end{column}
    \begin{column}{0.24\textwidth}
      \centering
      \begin{tikzpicture}[-, level distance=0.9cm,
          level 1/.style={sibling distance=1.3cm}]
        \node [arn_b, label=right:{\scriptsize fb $+2$}] {30}
        child { node [arn_b] {20}
            child { node [arn_b] {10} }
            child [missing]
          }
        child [missing];
      \end{tikzpicture}
    \end{column}
    \begin{column}{0.13\textwidth}
      \centering \scriptsize à direita\\ no 30
    \end{column}
    \begin{column}{0.24\textwidth}
      \centering
      \begin{tikzpicture}[-, level distance=0.9cm,
          level 1/.style={sibling distance=1.5cm}]
        \node [arn_b, label=right:{\scriptsize fb $0$}] {20}
        child { node [arn_b] {10} }
        child { node [arn_b] {30} };
      \end{tikzpicture}
    \end{column}
  \end{columns}

  \vspace{0.6cm}
  \centering
  O 30 tem fb $+2$ e o filho 10 tem fb $-1$, sinais contrários: é o
  zigue-zague. Rotação à esquerda no filho, que o alinha, e rotação à direita no
  nó.
\end{frame}

\begin{frame}[fragile]{Rebalancear}
\begin{minted}[fontsize=\scriptsize]{go}
func rebalanceia(a *No) *No {
	fb := fator(a)

	if fb > 1 { // pesada à esquerda
		if fator(a.Esq) < 0 { // caso esquerda-direita
			a.Esq = RotacaoEsquerda(a.Esq)
		}
		return RotacaoDireita(a)
	}

	if fb < -1 { // pesada à direita
		if fator(a.Dir) > 0 { // caso direita-esquerda
			a.Dir = RotacaoDireita(a.Dir)
		}
		return RotacaoEsquerda(a)
	}

	return a
}
\end{minted}
\end{frame}

\begin{frame}[fragile]{Inserção em AVL}
\begin{minted}[fontsize=\scriptsize]{go}
func Insere(a *No, v int) *No {
	if a == nil {
		return &No{Info: v}
	}

	if v < a.Info {
		a.Esq = Insere(a.Esq, v)
	} else if v > a.Info {
		a.Dir = Insere(a.Dir, v)
	} else {
		return a
	}

	atualizaAltura(a)
	return rebalanceia(a)
}
\end{minted}

  \vspace{0.2cm}
  A inserção da ABB com duas linhas a mais, executadas na volta da recursão em
  cada nó do caminho da folha nova até a raiz.
\end{frame}

\begin{frame}{Uma rotação por inserção}
  Uma inserção estraga o balanceamento de no máximo um nó. A rotação, simples ou
  dupla, devolve à subárvore a altura que ela tinha \textbf{antes} da inserção.

  \vspace{0.3cm}
  Como a altura da subárvore não mudou, nenhum nó acima dela precisa ser tocado:
  a inserção em AVL faz \tcr{no máximo uma} rotação.

  \vspace{0.5cm}
  A remoção não tem essa propriedade. Ela pode exigir rotações em todos os nós
  do caminho até a raiz, e é escrita do mesmo jeito: a remoção da ABB, seguida
  de \texttt{atualizaAltura} e \texttt{rebalanceia}.
\end{frame}

% ==========================================
\section{Custos}

\begin{frame}[fragile]{Entrada ordenada nas duas estruturas}
\begin{minted}{text}
Inserindo 1, 2, 3, ..., n em ordem crescente:
       n | altura na ABB | altura na AVL
---------+---------------+--------------
      10 |             9 |             3
     100 |            99 |             6
    1000 |           999 |             9
   10000 |          9999 |            13
\end{minted}

  \vspace{0.3cm}
  Com 10 mil chaves em ordem crescente, a busca na ABB percorre até 10 mil nós,
  e na AVL até 14. A altura 13 é a menor possível para 10 mil nós.
\end{frame}

\begin{frame}{ABB e AVL}
  \begin{center}
    \footnotesize
    \begin{tabular}{lll}
      \toprule
       & \textbf{ABB} & \textbf{AVL} \\
      \midrule
      Operações no caso médio & $\Theta(\log n)$ & $\Theta(\log n)$ \\
      Operações no pior caso & $\Theta(n)$ & $\Theta(\log n)$ \\
      Memória por nó & chave e 2 ponteiros & chave, 2 ponteiros e altura \\
      Rotações por inserção & nenhuma & no máximo 1 \\
      Rotações por remoção & nenhuma & até $O(\log n)$ \\
      \bottomrule
    \end{tabular}
  \end{center}

  \vspace{0.5cm}
  As árvores rubro-negras, na segunda metade da disciplina, admitem árvores um
  pouco mais desbalanceadas e fazem menos rotações na remoção, o que as torna a
  escolha mais comum em bibliotecas de uso geral. O critério de decisão é a
  proporção entre buscas e alterações na aplicação.
\end{frame}

\begin{frame}{De volta à medição da abertura}
  O slice ordenado é rápido na busca porque as chaves estão em ordem, e lento na
  inserção porque manter a ordem exige mover as chaves.

  \vspace{0.3cm}
  A árvore de busca guarda a mesma informação de ordem na forma da árvore, e por
  isso inserir custa descer até uma folha e pendurar um nó.

  \vspace{0.5cm}
  Aquela medição usou chaves em ordem aleatória, que é o caso médio. Com as
  mesmas 200 mil chaves em ordem crescente, a ABB daria uma coluna de 200 mil
  níveis, e as duas medições ficariam piores que as do slice.

  \vspace{0.3cm}
  A AVL remove a dependência da ordem de chegada.
\end{frame}

\begin{frame}{Exercícios}
  Os exercícios estão no roteiro da aula, na página do material. Eles
  \textbf{não valem nota} e não têm entrega: preparam a aula prática de
  árvores, essa sim avaliada, e o conteúdo cai na Prova 1.

  \vspace{0.4cm}
  Códigos da aula, para executar com \texttt{go run}:

  \begin{itemize}
    \item \texttt{06/codes/motivacao.go}: a medição da abertura;
    \item \texttt{06/codes/abb.go}: busca, inserção, mínimo e máximo;
    \item \texttt{06/codes/remocao.go}: os três casos da remoção;
    \item \texttt{06/codes/forma.go}: ordem de inserção e altura;
    \item \texttt{06/codes/avl.go}: rotações, os quatro casos e a comparação
          com a ABB.
  \end{itemize}
\end{frame}

\end{document}
